\documentclass[../main.tex]{subfiles}

\begin{document}
\subsection{Versión Compleja de las ecuaciones de Cauchy-Riemann}
\classdate{2026-04-13}

Sea \(F \colon U \subseteq \mathbb{C} \to \mathbb{C}\) con \(f = u + iv\). Decimos que
\(f\) cumple con las ecuaciones de Cauchy-Riemann en su forma compleja si

\begin{equation}
    \pdv{f}{x} = \pdv{u}{x} + i \pdv{v}{x}\qquad \wedge \qquad \pdv{f}{y} = \pdv{u}{y} + i
    \pdv{v}{y}
\end{equation}

existen y además

\begin{equation}
    \pdv{u}{y} + i \pdv{v}{x} = \pdv{f}{x} = i \pdv{f}{y} = - \pdv{v}{y} + i \pdv{u}{y}.
\end{equation}

Calculando \(\Re\) y \(\Im\) se enuncian de la manera tradicional las ecuaciones de
Cauchy-Riemann.
\end{document}
